\noindent{\heiti\zihao{4} 摘\quad 要：}\!以Quanser五杆绘图机器人的机构拓扑与笔式绘迹应用为参照，建立竖直板面椭圆轮廓描绘的课程设计模型，分析串联支链与并联整机的运动及负载关系。首先采用DH方法推导2R串联支链的正逆运动学与位置雅可比，基于动能和势能建立拉格朗日动力学，并用独立牛顿—欧拉递推进行核验；随后推导五杆机构的闭环正逆解、速度约束和奇异性条件，通过虚功将支链动力学投影为两个主动关节的整机动力学。Python计算表明，500组随机状态下两种串联动力学实现的最大扭矩差为$1.78\times10^{-15}\units{N\,m}$，高精度正动力学回放的末端最大偏差为$3.42\times10^{-8}\units m$。常规轨迹与近奇异轨迹的并联雅可比最大条件数分别约为5.56和273；在相同$1\units N$竖直等效力下，两个代表位形的单电机力矩幅值分别为0.896和$41.54\units{N\,m}$。结果说明，几何可达并不保证良好的传力性能，轮廓位置选择必须结合闭环奇异性评价。本文结果限于设定参数下的理想刚体仿真，可为绘图机构的轨迹布置与驱动负载分析提供可复现的课程算例。

\vspace{6pt}
\noindent{\heiti\zihao{4} 关键词：}\!五杆绘图机器人；串联与并联；雅可比矩阵；拉格朗日方程；牛顿—欧拉方程；奇异性
\vspace{5pt}
